\section{Related Work}
\label{sec:related}

We reviewed 26 studies in seven themes, each corresponding to a capability or a guarantee the system must provide. We summarise the themes and then state the relationship to the two closest systems.

\paragraph{Task interruption and resumption.}
A field study of office computer use established that people resume interrupted tasks from visible cues such as a window left open, and stopped at diagnosis~\cite{iqbal2007disruption}. ScreenTrack converted the finding into a tool by recording a time-lapse of the screen for visual scrolling, and its authors reported that retrieval slows as the history lengthens and that the approach suits visually distinct tasks~\cite{hu2020screentrack}. Our resumption feature records no screen; it works from the file index and the user's own notes, and because a workspace is bounded there is only a short list of recent items to present.

\paragraph{Bounded agents and trust.}
A trace-based assurance framework attaches machine-checkable contracts to each step an agent takes and mediates actions as allow, rewrite or block, but offers no way to undo an action once run~\cite{paduraru2026tracebased}. A survey of 51 attacks and 60 defences on agentic systems concludes that privilege separation and architectural confinement are more dependable than detecting malicious input, and documents that human approval degrades through decision fatigue and that most users cannot judge what they are asked to approve~\cite{kim2026attackdefense}. Both are written for organisations with security staff. Our workspace model is a plain-language rendering of the same principles for one person on a personal machine, and it adds reversal, which the assurance framework lacks. It does not solve approval fatigue, and we return to this in Section~\ref{sec:discussion}.

\paragraph{Agentic planning.}
Reflexion~\cite{shinn2023reflexion}, Toolformer~\cite{schick2023toolformer}, Generative Agents~\cite{park2023generative}, ReAct~\cite{yao2023react} and Tree of Thoughts~\cite{yao2023tot} establish reasoning, tool use and reflective memory as a workable substrate for agent behaviour. Each is evaluated on its own history, a sandbox world or benchmark tasks; none maintains a bounded workspace for a real user across interrupted work, and none makes consequential actions reversible. The failure modes they report, fabricated memory and repetitive loops, argue for keeping the number of unsupervised steps between user checkpoints small, which is how our planner is bounded.

\paragraph{Voice interaction.}
Two systematic reviews find voice usability studied in fragments and standard metrics silent on trust and cognitive load~\cite{deshmukh2024vui,dutsinma2022vareview}. Empirical work supplies design rules we adopt: an explanatory spoken indicator during long waits~\cite{speech2025waiting}, natural fillers rather than tones during delay~\cite{fillers2025mitigating}, and speech features as a real-time experience signal~\cite{ma2026beyondwords}. Voice as a complete mode over practical multi-step tasks remains under-explored, which is why our design makes it structurally identical to typing.

\paragraph{Code assistance with review.}
HULA deployed a plan, generate and approve loop at Atlassian and reported an 82 percent plan approval rate~\cite{takerngsaksiri2025hula}. HiLDe warns that regeneration can silently overwrite earlier edits and that too many alternatives overwhelm users~\cite{gonzalez2025hilde}; a study of suggestion timing warns against optimising acceptance rate~\cite{mozannar2024whento}; a taxonomy of expert steering strategies notes that some presuppose expertise~\cite{chidambaram2024sohf}; and an evaluation of models as code reviewers finds systematic overcorrection that longer explanations worsen~\cite{jin2026reliable}. Our code assistant therefore emits one proposed difference at a time and never writes to a file.

\paragraph{Faithfulness and citation.}
Hallucination appears in every abstractive summariser tested~\cite{maynez2020faithfulness}; FactCC is the most adoptable consistency checker~\cite{kryscinski2020evaluating}; no single metric detects every error category~\cite{pagnoni2021understanding}; and even frontier models judge faithfulness poorly~\cite{laban2023summedits}, which rules out asking the model to check its own citations. For citation, attribution has a usable definition~\cite{rashkin2023measuring}, incomplete citation is common in purpose-built systems~\cite{gao2023enabling}, and deployed generative search engines fail about a quarter of citations~\cite{liu2023evaluating}. Our synthesis operates over a small user-selected set and carries a citation map from each claim back to a chunk, which narrows the exposure, but the design contains no faithfulness check and does not verify that a citation supports its claim. These are open gaps rather than solved problems.

\paragraph{Closest systems.}
HULA and the agentic-security survey together occupy the two axes our work sits on: the approval loop as an interaction model, and confinement as a security position. We adopt both and claim neither. The differences are that scope is granted by the user rather than assumed, destruction is unavailable rather than guarded, reversal exists, approval is graduated by risk tier rather than uniform, the system discloses what it holds, the user is not assumed to be an engineer, and code assistance is one capability among several. Among deployed applications, inline coding assistants address code alone with whatever access the editor has, agentic command-line assistants execute across a machine and are hard to reverse, and conversational assistants hold no persistent relationship to the filesystem. Our formative study found that an experienced audience did not see the difference unaided, so the paper states it explicitly.
